 \section*{Resumen}

\begin{abstract}
Se desarrolló el modelo sísmico en ETABS de un edificio residencial de 21 pisos ubicado en Antofagasta, verificando el cumplimiento de los requisitos de la norma NCh433. El peso sísmico resultó en $P = 7{,}660$ tonf, un 12{,}6\% inferior al valor de prediseño, diferencia atribuible a las simplificaciones propias de esa etapa. Los períodos traslacionales fundamentales obtenidos son $T_X = 1{,}141$ s y $T_Y = 0{,}854$ s, coherentes con el período de prediseño de $T = 1{,}05$ s. Dado que los cortes basales iniciales no alcanzaban el mínimo normativo, los factores de reducción fueron ajustados a $R^*_X = 2{,}22$ y $R^*_Y = 2{,}69$, obteniéndose cortes basales de 746{,}41 tonf y 753{,}54 tonf en X e Y respectivamente, con una diferencia inferior al 3\% respecto del valor objetivo. Las derivas de entrepiso cumplen el límite admisible $\theta \leq 0{,}002$ en ambas direcciones. La rigidez lateral es adecuada en la dirección Y ($H/T = 64{,}22$ m/s), mientras que en X ($H/T = 48{,}05$ m/s) queda levemente fuera del rango esperado, en concordancia con su mayor período fundamental. Finalmente, el edificio no satisface completamente el criterio de desacoplamiento modal: la relación $T_\text{tors}/T_X = 1{,}10$ no supera el límite de 1{,}2, lo que sugiere que la respuesta sísmica en la dirección X puede verse influenciada por efectos torsionales.
\end{abstract}